The most common objection to algorithmic redistricting from state administrators is not technical---it is fiscal. ``How much will this cost?'' is the first question from legislative budget offices and redistricting staff. The answer is important: if algorithmic redistricting is expensive or administratively complex, its cost-benefit case is much weaker than if it is cheap and straightforward.

The answer is that algorithmic redistricting using the \texttt{bisect} platform is extraordinarily inexpensive by comparison to any alternative. The algorithm runs on public data that costs nothing to obtain, on commodity computing hardware, and produces a verifiable output in minutes. The dominant costs are human---staff time for quality assurance, public engagement, legal review, and documentation---and these costs are substantially lower than the human costs of any alternative redistricting process.

This paper is addressed to state administrators, legislative staff, budget officers, and procurement specialists who need to understand what it actually costs to implement algorithmic redistricting, what administrative law requires in terms of record-keeping and procedural compliance, and how to write a procurement specification that accurately describes the services they need.

Section~2 provides a full cost accounting for a \texttt{bisect} run, breaking costs into data acquisition, compute, staff time, audit, and legal review components. Section~3 compares these costs against the Iowa model (the cheapest existing alternative) and the California model (the most expensive). Section~4 analyzes APA record-keeping requirements and how the \texttt{bisect} manifest satisfies them. Section~5 provides a model procurement specification. Section~6 concludes.

One important framing note: the relevant comparison is not ``algorithmic redistricting cost'' versus ``zero,'' but ``algorithmic redistricting cost'' versus ``redistricting litigation cost.'' States that have litigated their redistricting maps have spent \$1M to \$40M in legal fees per cycle, with no guarantee of a neutral outcome. The question is not whether algorithmic redistricting costs money---it does---but whether it costs more or less than the alternatives.
